\documentclass[aspectratio=169,11pt]{beamer}

\usetheme{Madrid}
\usecolortheme{default}
\setbeamertemplate{navigation symbols}{}

\usepackage[T1]{fontenc}
\usepackage[utf8]{inputenc}
\usepackage{lmodern}
\usepackage{microtype}
\usepackage{hyperref}
\usepackage{xcolor}
\usepackage{ragged2e}

\definecolor{DeepNavy}{HTML}{1F2A44}
\definecolor{WarmGold}{HTML}{C58A2B}
\definecolor{SoftCream}{HTML}{F7F3EA}
\definecolor{SlateTeal}{HTML}{2F6E73}
\definecolor{SoftRed}{HTML}{8E2E2E}

\setbeamercolor{palette primary}{bg=DeepNavy, fg=white}
\setbeamercolor{palette secondary}{bg=SlateTeal, fg=white}
\setbeamercolor{palette tertiary}{bg=WarmGold, fg=white}
\setbeamercolor{structure}{fg=DeepNavy}
\setbeamercolor{frametitle}{bg=SoftCream, fg=DeepNavy}
\setbeamercolor{block title}{bg=DeepNavy, fg=white}
\setbeamercolor{block body}{bg=SoftCream, fg=black}
\setbeamercolor{alerted text}{fg=SoftRed}

\title{Part 3 -- Reflection Presentation}
\subtitle{Innovative Unit Plan: Minecraft Education Coding FUNdamentals}
\author{Piter Garcia}
\date{ED452B --- Spring 2026}

\begin{document}

\begin{frame}
  \titlepage
\end{frame}

\begin{frame}{What I designed}
\justifying
\begin{itemize}
  \item My innovative unit plan is a \textbf{4-lesson Grade 4 inclusive computer science unit} for Pine Brook's IGNITE rotation.
  \item I used \textbf{Minecraft Education Coding FUNdamentals} to teach sequencing, debugging, and algorithmic thinking.
  \item In placement, I kept seeing students get stuck on things like \textbf{menu clicks, sign-in, reading NPC directions, and managing multi-step agent code}.
  \item My main design goal was to make sure students were not labeled as \alert{``bad at coding''} when the real barrier might be reading load, navigation, language, or executive functioning.
  \item So the unit was built to separate \textbf{access barriers} from \textbf{actual computational thinking}. 
\end{itemize}
\end{frame}

\begin{frame}{Strategies I used in the unit plan}
\begin{itemize}
  \item \textbf{Started with placement evidence:} I used Pine Brook observations to identify real barriers students face in digital lessons.
  \item \textbf{Designed Tier 1 first:} predictable routine, explicit modeling, visual supports, partner structure, and low-stakes debugging norms.
  \item \textbf{Made the goal visible:} many students were stuck because they did not understand that \textbf{both gold plates had to be activated at the same time}.
  \item \textbf{Added targeted scaffolds:} chunking directions, vocabulary support, counting steps with students, and proximity support when needed.
  \item \textbf{Used multi-source assessment:} planning sheets, teacher observation, student explanation, completed code, and reflection --- not just task completion.
\end{itemize}
\end{frame}

\begin{frame}{Learning that shaped this plan}
\begin{itemize}
  \item \textbf{Karten and MTSS:} inclusion is not just placement; it is designing support into instruction from the start.
  \item \textbf{UDL thinking:} students need multiple ways to access content, show learning, and stay engaged.
  \item \textbf{Assessment course learning:} one score or one product never tells the full story of what a student knows.
  \item Placement made this real: some students needed a specific hint like ``your agent can float'' before their thinking became visible.
  \item \textbf{My own lens:} as a neurodivergent learner, I wanted to build a unit that does not confuse speed, compliance, or polished output with intelligence.
\end{itemize}

\begin{block}{Key shift in my thinking}
Inclusive planning is not adding accommodations at the end; it is building a lesson so more students can show what they know from the beginning.
\end{block}
\end{frame}

\begin{frame}{Reflection on the assignment}
\begin{itemize}
  \item This assignment helped me move from \textbf{general beliefs about inclusion} to \textbf{specific instructional decisions}.
  \item It made me think more carefully about how a lesson can look engaging on paper but still hide barriers for neurodivergent and CLD learners.
  \item It pushed me to connect theory to practice: Pine Brook observations, course readings, and my own teaching values all had to work together.
  \item The biggest lesson was that students often need support with the \textbf{entry into the task}, not just the concept itself.
  \item If I revise this unit again, I want to use even more implementation data to see which supports truly improved access and which need adjustment.
\end{itemize}
\end{frame}

\begin{frame}{What I am taking forward}
\begin{itemize}
  \item Start with \textbf{learner variability}, not the myth of an average student.
  \item Treat assessment as a way to \textbf{surface strengths and thinking}, not just sort students.
  \item Build supports early so they feel normal, not like labels.
  \item Keep asking: \textbf{Is this task measuring the learning goal, or just the student's ability to get around the barriers I created?}
\end{itemize}

\vfill
\centering
\textbf{Thank you}
\end{frame}

\end{document}
